\documentclass[fleqn,usenatbib]{mnras}

\usepackage{amsmath}
\usepackage{amssymb}
\usepackage{booktabs}
\usepackage{graphicx}
\usepackage{hyperref}
\usepackage{xcolor}

\newcommand{\ROBERT}{\textsc{Robert}}
\newcommand{\roberttodo}[1]{\textcolor{magenta}{[TODO: #1]}}
\newcommand{\code}[1]{\texttt{#1}}

\title[ROBERT]{ROBERT: a reproducible retrieval framework for JWST exoplanet emission spectra}

\author[J. Taylor et al.]{
Jake Taylor$^{1}$\thanks{E-mail: \roberttodo{corresponding-author email}}
and \roberttodo{co-authors}\\
$^{1}$\roberttodo{affiliation}
}

\date{Accepted XXX. Received YYY; in original form ZZZ}
\pubyear{2026}

\begin{document}
\label{firstpage}
\pagerange{\pageref{firstpage}--\pageref{lastpage}}
\maketitle

\begin{abstract}
Atmospheric retrievals translate exoplanet spectra into constraints on thermal
structure, composition, clouds, and planetary properties, but scientific
interpretation can depend on choices made across opacity preparation,
radiative transfer, inference, and data handling. We present \ROBERT{}, an
open, typed Python framework developed for reproducible retrievals of JWST
exoplanet emission spectra. \ROBERT{} separates atmospheric state construction,
correlated-$k$ opacity preparation, radiative transfer, instrument response,
likelihood evaluation, and inference behind explicit interfaces, and records
the resulting configuration and data provenance. The current implementation
supports clear and cloudy thermal emission, transmission, multiple one-dimensional
regions, multiple data sets, optimal estimation, and nested sampling with
UltraNest. We validate individual numerical components and compare thermal
radiative transfer with PICASO and molecular emission and transmission with
petitRADTRANS. We further exercise the end-to-end data model using the
panchromatic JWST emission spectrum of WASP-69b. \roberttodo{Replace this final
sentence with the quantitative WASP-80b results, including retrieved metallicity,
C/O, molecular detections, and model evidences.} This work establishes a
testable route from retrieval-code development to a reproducible scientific
analysis while retaining clear boundaries between validated capability and
future physics.
\end{abstract}

\begin{keywords}
planets and satellites: atmospheres -- planets and satellites: gaseous planets
-- methods: numerical -- radiative transfer -- software: development
\end{keywords}

\section{Introduction}
\label{sec:introduction}

Atmospheric retrieval is the inverse problem of inferring atmospheric and
planetary properties from a measured spectrum. A forward model maps a parameter
vector $\boldsymbol{\theta}$ to a spectrum $\boldsymbol{m}(\boldsymbol{\theta})$;
an inference method then evaluates which regions of parameter space are
consistent with the observation. The method has developed from terrestrial and
Solar-system remote sensing, including optimal estimation and the NEMESIS
framework \citep{Rodgers2000,Irwin2008}, into a central tool for interpreting
transiting exoplanets \citep[e.g.][]{MadhusudhanSeager2009,Line2013}.

Modern retrieval packages make different trade-offs between physical breadth,
speed, configurability, and software transparency. TauREx provides a modular
Python retrieval architecture and a plugin system \citep{AlRefaie2021};
petitRADTRANS combines line-by-line and correlated-$k$ radiative transfer with
retrieval tooling \citep{Molliere2019,Molliere2020}; POSEIDON emphasizes
multidimensional atmospheres and broad observing geometries
\citep{MacDonald2017,MacDonald2024}; and PICASO couples reflected-light,
thermal-emission, transit, climate, and scattering calculations
\citep{Batalha2019,Batalha2022}. CHIMERA, Brewster, ARCiS, PLATON, and
NEMESIS/NemesisPy span complementary use cases from hot-Jupiter transmission
and emission to directly imaged planets and brown dwarfs
\citep{Line2013,Burningham2017,Min2020,Zhang2019,Irwin2008}.

No single framework is uniformly optimal. Broad packages can make many science
cases accessible, but tightly coupled data loading, physics, and inference can
make controlled replacement or validation of one component difficult.
Research-oriented codes may be fast and scientifically productive while
depending on project-specific scripts, implicit conventions, or external data
layouts. Conversely, strict abstraction can obscure the numerical path or
introduce complexity before the underlying physics is validated. These are
engineering trade-offs rather than a ranking of scientific quality.

JWST has made these trade-offs especially consequential. Its precision and
broad wavelength coverage expose molecular features, cloud effects, and
longitudinal temperature structure simultaneously. For example, the
$2$--$12\,\mu$m emission spectrum of WASP-69b is inconsistent with a simple
homogeneous clear atmosphere and favours models containing aerosols and a
dayside temperature gradient \citep{Schlawin2024}. WASP-80b provides a cooler
comparison target: NIRCam transmission and emission spectra yielded strong
methane evidence \citep{Bell2023}, and its later panchromatic emission spectrum
constrained metallicity and C/O over $2.4$--$12\,\mu$m
\citep{WASP80Panchromatic2025}.

Here we introduce \ROBERT{} (\roberttodo{expand the acronym, or state explicitly
that ROBERT is a name rather than an acronym}), a Python package built around
explicit scientific data models and independently testable numerical
components. Its initial science focus is JWST thermal emission, while its
interfaces avoid encoding an emission-only assumption. The distribution is
named \code{robert-exoplanets} and is currently pre-release software.

This paper has two aims. First, we describe the architecture, physics, inference
interfaces, and validation strategy of \ROBERT{}. Secondly, we use a reanalysis
of WASP-80b to test whether a reproducible end-to-end implementation recovers
the published atmospheric constraints and whether additional temperature,
cloud, chemistry, or multidataset flexibility is supported by the data. Section
\ref{sec:methods} describes the model and benchmarks, Section
\ref{sec:results} defines and will report the WASP-80b analysis, and Section
\ref{sec:conclusions} summarizes the present scope and future development.


\section{Methods}
\label{sec:methods}

\subsection{Design and scientific data flow}

\ROBERT{} separates six stages: observation ingestion, atmosphere construction,
opacity evaluation, radiative transfer, instrument response, and inference.
Typed immutable containers carry pressure and spectral grids, stellar and
planetary properties, atmospheric state, observations, and spectra. A retrieval
parameter vector is therefore distinct from the physical atmosphere that it
generates. This separation is intended to make conventions---including units,
grid orientation, reference radius, opacity identity, and interpolation
policy---visible at component boundaries.

For data set $d$, the forward path is
\begin{equation}
 \boldsymbol{\theta}\rightarrow
 \mathcal{A}(\boldsymbol{\theta})\rightarrow
 \mathcal{K}\rightarrow
 \mathcal{R}\rightarrow
 \mathcal{I}_d\rightarrow
 \boldsymbol{m}_d(\boldsymbol{\theta}),
 \label{eq:dataflow}
\end{equation}
where $\mathcal{A}$ constructs the atmosphere, $\mathcal{K}$ evaluates opacity,
$\mathcal{R}$ solves radiative transfer, and $\mathcal{I}_d$ applies the data
set's response and binning. Opacity preparation is cached outside likelihood
calls. Run manifests record the model configuration, priors, opacity and input
checksums, software environment, random seed, and grid hashes.

\subsection{Atmospheric structure and composition}

The atmosphere is represented on an explicitly ordered pressure grid. Available
temperature prescriptions include isothermal, tabulated, Guillot-type
\citep{Guillot2010}, Madhusudhan--Seager-type \citep{MadhusudhanSeager2009},
and Parmentier--Guillot-type profiles \citep{ParmentierGuillot2014}.
Composition may be supplied as free constant-with-altitude mixing ratios,
layer-dependent profiles, or equilibrium chemistry through FastChem
\citep{Stock2018}. Trace-gas abundances are combined with an explicit H$_2$/He
background, with validation of non-negative abundances and normalization.

Under hydrostatic equilibrium, the absorber column in layer $i$ is computed as
\begin{equation}
 N_i = \frac{\Delta P_i}{\bar{\mu}_i m_{\rm u} g_i},
 \label{eq:column}
\end{equation}
where $\Delta P_i$ is the pressure difference, $\bar{\mu}_i$ the mean molecular
weight in atomic mass units, $m_{\rm u}$ the atomic mass constant, and $g_i$
the gravitational acceleration. Spherical shell radii can be obtained by
integrating hydrostatic balance from a declared reference pressure and radius.

\subsection{Opacity and correlated-$k$ treatment}

\ROBERT{} reads NEMESIS KTA tables and opacity archives prepared with
\code{exo\_k} \citep{Leconte2021}. Gas opacity is interpolated in pressure and
temperature under an explicit bounds policy and evaluated on the model's
spectral and $g$ grids. For gas $s$, layer $i$, wavelength bin $j$, and ordinate
$l$, the monochromatic-distribution optical depth is
\begin{equation}
 \tau_{sijl}=k_{sijl}\,x_{si}\,N_i,
 \label{eq:tau}
\end{equation}
with unit conversion determined by opacity metadata. Multiple gases may share
an ordinate grid or be combined by random overlap with resorting and
rebinning. Collision-induced absorption and Rayleigh extinction are separate
contributors, allowing their provenance and approximations to be inspected.

For random overlap, two discrete opacity distributions with optical depths
$\{\tau_a\}$ and $\{\tau_b\}$ and quadrature weights $\{w_a\}$ and $\{w_b\}$
form the product distribution
\begin{equation}
 \tau_{ab}=\tau_a+\tau_b, \qquad w_{ab}=w_aw_b.
 \label{eq:random-overlap}
\end{equation}
The $n_g^2$ product points are ranked and conservatively rebinned to the
original $n_g$ intervals; the procedure is applied successively to all active
absorbers. The reference implementation explicitly sorts the product
distribution. In the accelerated implementation, nondecreasing input
distributions---the normal correlated-$k$ case---are recognized and the
$n_g$ already sorted rows of outer sums are combined with an $n_g$-element
heap. Rebinning is performed in the same streaming pass with reusable work
arrays. This avoids a general sort of every $n_g^2$ temporary at every layer
and wavelength. Unsorted inputs retain the general sorting path, and a gas is
skipped only when its maximum optical depth is below the declared transparent
gas threshold. The NumPy implementation remains the scientific reference;
the optional compiled backend is required to reproduce it to floating-point
precision.

\subsection{Profile-guided computational optimization}
\label{sec:performance-methods}

Optimization used warmed forward-model profiles of an 80-layer, six-molecule
WASP-69b calculation with 16 $g$ ordinates, evaluated on both a 1591-point
native opacity grid and 280-point observation grids. Changes were accepted
only if they preserved the layer, wavelength, angular, and correlated-$k$
resolution, equations, convergence thresholds, and double-precision
arithmetic. Consequently, the changes described here alter implementation
cost rather than retrieval physics.

Random overlap was the dominant initial cost and was optimized as described
above. Two additional repeated operations were removed from likelihood calls.
First, prepared opacity objects retain validated spectral lookup indices, and
providers may cache $\log k$ for immutable coefficient tables instead of
recomputing logarithms during every pressure--temperature interpolation. The
cache is explicit and optional because it exchanges memory for speed; cached
and uncached calculations are tested for exact equality. Second, four fixed
four-term source contractions in the SH4 solver were written as explicit sums,
avoiding thousands of dispatches to a general tensor-contraction routine while
leaving the spherical-harmonics equations unchanged. Spherical transmission
was not rewritten because profiling showed that its contribution was already
small.

The random-overlap benchmark spans one, three, and six species at 280 and 1591
wavelengths, records throughput and memory, and compares a deterministic
subset with the NumPy reference on every run. On the development system
(macOS arm64, Python 3.12.13, NumPy 2.2.6, and Numba 0.61.2), the largest
observed random-overlap difference was $3.13\times10^{-13}$ in optical depth.
The complete 279-test suite passes after these changes. In the representative
WASP-69b workloads, the native-grid-then-bin path decreased from 6.35 to
1.21\,s per warmed evaluation (a factor of 5.3), while preparing opacities per
instrument mode decreased from 1.34 to 0.43\,s (a factor of 3.1). The
corresponding log-likelihoods were unchanged exactly within the recorded
double-precision calculation. Cached log-coefficients used approximately
106\,MiB for the four observation-grid providers and 599\,MiB for the full
native-grid provider, so production configurations should select the cache
policy with available memory in mind.

\subsection{Thermal emission and transmission}

For a non-scattering, plane-parallel ray with direction cosine $\mu$, the
top-of-atmosphere intensity is evaluated from the formal solution
\begin{equation}
 I_{\lambda}(0,\mu) = I_{\lambda}(\tau_{\rm b},\mu)
 e^{-\tau_{\rm b}/\mu} +
 \int_{0}^{\tau_{\rm b}} B_{\lambda}[T(t)]
 e^{-t/\mu}\frac{\mathrm{d}t}{\mu}.
 \label{eq:formal}
\end{equation}
When temperatures are specified at layer edges, \ROBERT{} integrates a Planck
source that is linear in optical depth within each layer; otherwise it records
the constant layer-centre approximation. Disc integration supports normal
emission, a uniform thermal disc, and Lobatto quadrature. Eclipse depths divide
the planet flux by either an input stellar spectrum or a blackbody stellar
reference.

Clouds carry extinction optical depth $\Delta\tau$, single-scattering albedo
$\omega_0$, and asymmetry factor $g$. A grey mass-extinction coefficient
$\kappa_{\rm cld}$ gives
\begin{equation}
 \Delta\tau_{{\rm cld},i}=0.1\,
 \kappa_{\rm cld}\frac{\Delta P_i}{g_i},
 \label{eq:grey-cloud}
\end{equation}
where $\kappa_{\rm cld}$ is in cm$^2$ g$^{-1}$ and the factor $0.1$ converts to
SI units. Multiple scattering can be solved with a hemispheric-mean two-stream
method based on \citet{Toon1989} or a four-term spherical-harmonics (SH4)
method following \citet{Rooney2023}, including Henyey--Greenstein moments and
optional delta-$M$ scaling. The current transmission path integrates slant
optical depths through hydrostatic spherical shells; \roberttodo{state the
validated refraction, scattering, and cloud limitations for the release used
in the paper}.

\subsection{Inhomogeneous daysides and multiple data sets}

The emergent spectrum of a two-region dayside is mixed after radiative transfer,
\begin{equation}
 F_{\lambda}=f_{\rm hot}F_{\lambda,\rm hot}+
 (1-f_{\rm hot})F_{\lambda,\rm cold},
 \label{eq:two-region}
\end{equation}
where the columns may have independent thermal and cloud structures while
sharing selected chemistry parameters. This ordering is essential because
radiative transfer and the Planck function are nonlinear. A diluted one-region
model is the limiting case in which the complementary region contributes
negligible flux. Each observation retains its instrument identity, wavelength
bins, mask, and nuisance parameters; the joint likelihood is the sum of the
per-data-set terms.

\subsection{Likelihoods and inference}

For independent Gaussian uncertainties, the log likelihood is
\begin{equation}
 \ln \mathcal{L}=-\frac{1}{2}\sum_d\sum_i
 \left[\frac{(y_{di}-m_{di})^2}{s_{di}^2}
 +\ln(2\pi s_{di}^2)\right],
 \quad s_{di}^2=\sigma_{di}^2+\sigma_{\rm jit,d}^2.
 \label{eq:likelihood}
\end{equation}
The same sampler-independent retrieval problem can be evaluated with optimal
estimation \citep{Rodgers2000} or nested sampling. The current Bayesian adapter
uses UltraNest's reactive nested sampler \citep{Buchner2021}; priors and their
unit-cube transforms are explicit objects. Model comparisons will report
evidence differences only after convergence and repeated-run checks.

\subsection{Numerical verification and benchmarking}
\label{sec:benchmarks}

Validation is layered. Unit tests cover grids, conservation and limiting cases,
opacity interpolation, optical-depth assembly, radiative-transfer limits,
likelihoods, manifests, and inference adapters. Injection--recovery tests then
exercise the complete parameter-to-spectrum path. Cross-code comparisons use
matched atmospheric states and, where possible, identical optical-depth arrays
to isolate the radiative-transfer kernel from opacity differences.

For a controlled 160-layer comparison, matched \ROBERT{} and PICASO SH4 point
radiances agree to a maximum relative difference of $3.23\times10^{-6}$ for
both isotropic scattering and a case with $\omega_0=0.9$ and $g=0.6$. Using
real molecular correlated-$k$ structure for HAT-P-32b, the maximum and RMS
disc-integrated differences are currently $0.0507$ and $0.0151$ per cent,
respectively; a FastChem plus collision-induced-absorption case reaches a
maximum difference of $0.00361$ per cent. \roberttodo{Regenerate these values
from a tagged release and add a table containing grid, hardware, package
versions, metric definitions, and output-file hashes.}

The petitRADTRANS benchmark suite compares stable reference calculations for
multispecies emission and transmission and tests convergence with native HDF5
opacities. \roberttodo{Insert the final petitRADTRANS version, exact common
inputs, residual metrics, and figure after freezing the reference data.}

\subsection{WASP-69b end-to-end benchmark}

The WASP-69b benchmark uses the 280-point published $2.455$--$11.875\,\mu$m
spectrum of \citet{Schlawin2024}, retaining NIRCam/F322W2, the NIRCam overlap,
NIRCam/F444W, and MIRI/LRS as four named data sets. We compare one-region clear,
two-region clear, diluted one-region, and two-region grey-cloud models. The
primary validation quantities are the recovered hot-region fraction, thermal
profiles, cloud contrast, metallicity, C/O, maximum likelihood, and evidence
relative to the published analysis. \roberttodo{Run and report the staged
retrieval suite; the current repository contains the data model and target
workflow but not final publication-grade posteriors.}

\section{WASP-80b reanalysis}
\label{sec:results}

This section defines the analysis that will provide the principal science
result of the paper. It intentionally contains no inferred values before the
retrieval and robustness suite is complete.

\subsection{Data and model hierarchy}

We will reanalyse the published NIRCam/F322W2, NIRCam/F444W, and MIRI/LRS
$2.4$--$12\,\mu$m dayside emission spectrum of WASP-80b
\citep{Bell2023,WASP80Panchromatic2025}. The reduction, wavelength bins,
uncertainties, covariance assumptions, and any inter-instrument offsets will be
recorded as separate observation data sets rather than concatenated into an
anonymous spectrum. \roberttodo{Acquire the exact machine-readable spectrum,
record its DOI/MAST identifiers and checksums, and determine whether the
published covariance is available.}

The planned hierarchy is:
\begin{enumerate}
 \item a clear one-region free-chemistry model;
 \item a clear one-region equilibrium-chemistry model parameterized by
 metallicity and C/O;
 \item grey and microphysical cloud extensions;
 \item diluted and two-region dayside models; and
 \item targeted disequilibrium or quench-pressure extensions if justified by
 posterior predictive residuals.
\end{enumerate}
Each comparison will hold opacity, spectral response, priors on shared
parameters, and convergence criteria fixed where scientifically possible.

\subsection{Prespecified reported quantities}

The primary quantities are metallicity, C/O, thermal structure, CH$_4$ and
H$_2$O abundance constraints, cloud optical properties, and the evidence of
each model relative to the clear one-region baseline. We will report credible
intervals, prior-to-posterior contraction, posterior correlations, contribution
functions, and posterior predictive spectra. Molecular evidence will be tested
with nested model comparisons or abundance-posterior diagnostics declared
before inspecting the final result; a Gaussian ``sigma'' equivalent will not
be quoted without stating its definition.

\subsection{Robustness tests}

Robustness tests will include NIRCam-only and MIRI-only retrievals, alternative
temperature parameterizations, opacity-table resolution and bounds checks,
cloud-model sensitivity, removal of instrument offsets, repeated nested-sampling
runs, and injection--recovery at representative posterior points. The
published WASP-80b posterior products will be passed through the same plotting
and comparison code where their parameterization permits.

\subsection{Results completion gate}

Before this section is converted from a plan into results, all final runs must
have: (i) a versioned configuration and manifest; (ii) converged evidence and
posterior diagnostics; (iii) reproducible figures and tables; (iv) an archived
machine-readable posterior; and (v) an internal check that the text distinguishes
data-supported conclusions from model-dependent interpretation.

\roberttodo{Replace this paragraph with the quantitative results and add: data
figure; best-fitting/posterior-predictive spectra; abundance and thermal-profile
figure; corner or information-content summary; evidence table; and direct
comparison with \citet{WASP80Panchromatic2025}.}


\section{Conclusions}
\label{sec:conclusions}

We have presented \ROBERT{}, a pre-release Python framework for testable and
reproducible exoplanet atmospheric retrievals. Its principal contribution is a
set of explicit boundaries between observations, atmospheric state, opacity,
radiative transfer, instrument response, likelihoods, and inference. The
current implementation includes correlated-$k$ gas opacity and random overlap,
thermal emission and transmission, clear and scattering-cloud calculations,
one- and two-region daysides, multiple observation data sets, optimal
estimation, and UltraNest.

Controlled comparisons with PICASO already test the SH4 multiple-scattering
path at high numerical precision, while petitRADTRANS comparisons and the
WASP-69b data workflow broaden validation to molecular opacity, transmission,
and real multi-instrument observations. \roberttodo{Update this sentence after
all tagged benchmark artefacts and WASP-69b retrievals are complete.}

The WASP-80b analysis is designed as both a scientific reanalysis and a test of
the complete reproducibility chain. \roberttodo{Insert the main WASP-80b
findings and state clearly where they agree or disagree with published
constraints.}

Near-term development will focus on science-grade cloud microphysics,
high-stream scattering references, broader opacity validation, richer
disequilibrium chemistry, and robust long-run posterior workflows. Longer-term
interfaces permit reflected light, phase curves, multidimensional atmospheres,
line-by-line opacity, additional samplers, and multi-instrument covariance,
without requiring those capabilities to be claimed by the present release.


\section{Software and reproducibility}
\label{sec:reproducibility}

The manuscript source is maintained with the \ROBERT{} code and is structured
for two-way manual synchronisation between GitHub and Overleaf. Analysis code,
configuration objects, tests, and manuscript figures will be versioned together.
Every publication run will write a manifest before inference begins and stable
machine-readable result files after completion. The archived release will pin
the \code{robert-exoplanets} version and dependencies, record opacity and data
checksums, and provide commands for reproducing every table and figure.

\roberttodo{Add the tagged software release, Zenodo DOI, licence, Python
versions, hardware used for timings, complete environment lock file, and a
table mapping each paper figure/table to its generating command and output.}



\section*{Acknowledgements}
\roberttodo{Funding, computing facilities, and contributors.} The development
and writing workflow used OpenAI Codex. The authors reviewed and take
responsibility for all code, analyses, citations, and text produced with AI
assistance.

\section*{Data Availability}
The WASP-69b spectrum used for validation is distributed with source provenance
in the \ROBERT{} repository and originates from the machine-readable tables of
\citet{Schlawin2024}. \roberttodo{Add the WASP-80b data products, MAST programme
identifier, exact reduction version, retrieval posterior archive, and Zenodo
DOI.} Source code and analysis scripts will be archived at a versioned DOI on
acceptance; the development repository is \roberttodo{insert public GitHub URL}.

\bibliographystyle{mnras}
\bibliography{references}

\label{lastpage}
\end{document}

